\documentclass[11pt]{article}
\usepackage[T1]{fontenc}
\usepackage{lmodern}
\usepackage[margin=1in]{geometry}
\usepackage{amsmath,amssymb,amsthm}
\usepackage{booktabs,array,longtable}
\usepackage{hyperref}
\hypersetup{hidelinks,pdftitle={A counterexample to packing-domination inheritance in strong products}}
\allowdisplaybreaks
\raggedbottom
\newtheorem{theorem}{Theorem}[section]
\newtheorem{lemma}[theorem]{Lemma}
\newtheorem{proposition}[theorem]{Proposition}
\newtheorem{corollary}[theorem]{Corollary}
\theoremstyle{definition}
\newtheorem{definition}[theorem]{Definition}
\theoremstyle{remark}
\newtheorem{remark}[theorem]{Remark}
\newcommand{\F}{\mathbb F}
\newcommand{\wt}{\operatorname{wt}}
\newcommand{\dist}{\operatorname{dist}}
\newcommand{\supp}{\operatorname{supp}}
\newcommand{\Reach}{\operatorname{Reach}}
\newcommand{\code}{\mathcal C}
\newcommand{\pairs}{\mathcal P_4}
\title{A counterexample to packing-domination\\
inheritance in strong products}
\author{River Zhang\\[3pt]\small Chengdu Neusoft University}
\date{Preprint draft v1.2 --- 10 October 2026}

\begin{document}
\maketitle
\begin{abstract}
The \(d\)-distance \(p\)-packing domination number of a graph is the
minimum size of a set that dominates within distance \(d\) and separates
distinct centers by distance at least \(p+1\), with value infinity when no
such set exists. Bujt\'as, Ir\v si\v c Chenoweth, Klav\v zar and Zhang
conjectured that nonexistence in one factor is inherited by every strong
product containing that factor. We construct a finite connected simple
graph \(H\), with 608 vertices and 9,168 edges, for which
\[
 \gamma_2^3(Q_6)=\infty
 \qquad\text{but}\qquad
 \gamma_2^3(Q_6\boxtimes H)\le 64.
\]
The auxiliary graph uses singleton and distance-four pair supports in
the binary six-cube. Compatibility of support unions prevents short
paths between incompatible code labels, while explicit eight-label
sets provide radius-two coverage at pair vertices. We give a complete
proof from actual graph adjacency. A Lean formalization checks the
complete existence-level counterexample on finite simple graphs and
standard graph walks; two separately implemented exact checkers
reconstruct and verify the finite graph and its same-center covering
witness. The numerical minimum is not claimed to be 64, and historical
priority is not asserted.
\end{abstract}
\noindent\textbf{Keywords.} Distance domination; packing; strong graph
product; hypercube; explicit counterexample; formal verification.
\par\noindent\textbf{2020 Mathematics Subject Classification.}
Primary 05C69; Secondary 05C76, 68V15.

\medskip
\noindent\textbf{Draft status.} The author name and affiliation were supplied
by the project owner as draft metadata; accountable human examination and publication approval remain
pending. This document has not been submitted to arXiv
or a journal and has not undergone human peer review. Substantial AI
participation is disclosed in Section~\ref{sec:disclosure}.
Version v1.2 retains the mathematical construction of v1, expands the
literature discussion, and supplies a versioned reproducibility package.

\section{Introduction}\label{sec:introduction}
Distance domination and packing impose opposing requirements on a set of
graph vertices. A dominating set must place centers near all vertices,
whereas a packing must keep distinct centers apart. Their simultaneous
feasibility is therefore a structural question, rather than merely an
optimization problem. Following Bujt\'as, Ir\v si\v c Chenoweth,
Klav\v zar and Zhang~\cite{bickz-strong}, we write
\(\gamma_d^p(F)\) for the minimum cardinality of a \(d\)-distance
dominating \(p\)-packing in a graph \(F\). The value is \(\infty\) when
there is no feasible set.
The underlying joint notion was introduced by Beineke and
Henning~\cite{beineke-henning} using the older
\((p,d)\)-domination notation; our notation follows the contemporary
source.

For the strong product \(G\boxtimes H\), the standard distance formula is
\[
 \dist_{G\boxtimes H}((g,h),(g',h'))
   =\max\{\dist_G(g,g'),\dist_H(h,h')\}.
\]
One can consequently combine feasible sets from both factors by taking
their Cartesian product. The converse existence question is less
immediate: can the second factor supply separation missing in the first
while still allowing the same product centers to dominate every layer?
Conjecture~3.1 of~\cite{bickz-strong} proposes the implication
\begin{equation}\label{eq:original-conjecture}
 \gamma_d^p(G)=\infty
 \quad\Longrightarrow\quad
 \gamma_d^p(G\boxtimes H)=\infty
 \quad\text{for every graph }H.
\end{equation}
The hypothesis concerns only \(G\); it does not require both factors
to have infinite packing-domination number.

The source proves the conjectured implication for several classes and
parameter ranges. It identifies \(d<p<2d\) as the remaining nontrivial
interval. When \(p=2d\), feasible sets are perfect \(d\)-codes: their
closed radius-\(d\) balls form a partition. Perfect-code inheritance in
strong products is established in~\cite{perfect}. Our example instead
uses \(d=2,p=3\), where balls associated with a packing may overlap.
This is the smallest integral parameter pair in the remaining interval.
The support construction below uses that overlap directly.
Thus the result concerns existence of a feasible set, before any
comparison of finite optimal values. A product can restore feasibility
even when one factor admits no feasible set at all.

\subsection{Main result}

\begin{theorem}\label{thm:main}
There is a finite nonempty connected simple undirected graph \(H\) with
608 vertices and 9,168 edges such that
\begin{equation}\label{eq:main}
 \gamma_2^3(Q_6)=\infty,
 \qquad
 \gamma_2^3(Q_6\boxtimes H)\le 64.
\end{equation}
More precisely, \(H\) is the graph of Definition~\ref{def:H}, and
\(\code=\{(a,A_a):a\in\F_2^6\}\) is a \(2\)-distance dominating
\(3\)-packing in its strong product with \(Q_6\).
\end{theorem}

The theorem contradicts~\eqref{eq:original-conjecture} with graphs that
are finite, simple and connected. The original conjecture does not
restrict \(H\) to a cycle, to a regular graph, or to a graph isomorphic
to \(G\). Our auxiliary factor is designed to couple the two
requirements. Its adjacency is explicit; no desired distance table
is postulated.

The proof separates into three tasks. First, a classical binary
distance-sum argument and a small structural calculation show that no
subset of \(Q_6\) is feasible. Second, compatibility of supports excludes
all paths of length at most three between code vertices whose labels
are too close in \(Q_6\). Third, each auxiliary vertex sees, within
two steps, a set of code labels covering the entire cube within
distance two. The final step uses one label for both coordinates;
separate covering witnesses in the two factors would not suffice.

All mathematical arguments needed for the theorem are given here.
Sections~\ref{sec:formal} and~\ref{sec:computational} describe additional
kernel and exact-program checks without replacing the proof.
The distance-sum method is classical Plotkin counting~\cite{plotkin},
and the product distance formula is standard~\cite{handbook}.
Compatibility graphs and clique supports are also established general
tools. We do not claim novelty for these ingredients or assert
historical firstness for this counterexample.

\section{Related work}\label{sec:related}
Beineke and Henning~\cite{beineke-henning} studied independent distance
domination under a joint condition with separation parameter \(p\)
and covering radius \(d\). The
parameter in this paper is the same simultaneous packing and domination
requirement, with notation adopted from~\cite{bickz-strong}. It must not
be confused with the two separate optimization parameters for packing
and domination. Fisher~\cite{fisher} studied domination, fractional
domination, 2-packing and graph products. We cite that work for the
older product context, without attributing to it an unverified
statement about the joint parameter considered here.

Abay-Asmerom, Hammack and Taylor~\cite{perfect} established the
inheritance of perfect \(r\)-code nonexistence in strong products.
Their result applies to the endpoint \(p=2d\), where the closed
dominating balls partition the vertex set. The conjecture
in~\cite{bickz-strong} extends this existence question to general
packing-domination parameters; our counterexample lies strictly
between \(d\) and \(2d\). It therefore preserves the perfect-code
result while disproving the proposed extension. The close-vertex
hypothesis used for a positive result in~\cite{bickz-strong} is also
absent from our auxiliary graph, as proved in
Proposition~\ref{prop:no-close}.

The companion work~\cite{bickz-complexity} investigates computational
complexity, exact values for paths and cycles, and bounds for trees.
Those results do not impose
such a restriction on the universal quantifier over \(H\) in
Conjecture~3.1. Our auxiliary graph need not be a cycle, a tree,
or a second hypercube. The source correspondence is stated explicitly
in~\eqref{eq:original-conjecture}, and all graphs used here are finite
and connected.

The counting argument for the cube obstruction is a classical
Plotkin-type method~\cite{plotkin}; the strong-product metric is
standard~\cite{handbook}. Neither ingredient is claimed as a new
theorem. The contribution presented here is the explicit support
construction and its complete verification as a counterexample to
the stated universal implication. Targeted literature checks have
not established historical firstness; their limitations are recorded
in Section~\ref{sec:discussion}. Lean~\cite{lean4} and
Mathlib~\cite{mathlib} provide an additional, separately delimited
formal check of the existence assertion, described in
Section~\ref{sec:formal}.

\section{Definitions and graph distances}\label{sec:preliminaries}
Throughout the proof, a graph is finite, nonempty, simple and undirected;
graphs whose distances we use are connected. A walk is allowed to
repeat vertices. The distance \(\dist_F(u,v)\) is the least length of
a \(u\)-to-\(v\) walk in \(F\). For \(r\ge0\), let
\[
 B_r^F(u)=\{v\in V(F):\dist_F(u,v)\le r\}.
\]
We suppress the superscript when the ambient graph is clear.

\begin{definition}\label{def:packing-domination}
Let \(d,p\) be nonnegative integers. A set \(D\subseteq V(F)\) is a
\emph{\(d\)-distance dominating set} if
\[
 \forall v\in V(F)\ \exists a\in D:\ \dist_F(v,a)\le d.
\]
It is a \emph{\(p\)-packing} if
\[
 a,b\in D,\ a\ne b\quad\Longrightarrow\quad \dist_F(a,b)\ge p+1.
\]
The packing-domination number is
\[
 \gamma_d^p(F)=
 \begin{cases}
 \min\{|D|:D\text{ has both properties}\},&
       \text{if a feasible set exists},\\
 \infty,&\text{otherwise}.
 \end{cases}
\]
\end{definition}

The source~\cite{bickz-strong} states the domination condition only for
vertices outside \(D\). This is equivalent to the closed-ball formulation
above, since a center dominates itself at distance zero and \(d\ge0\).
For a finite graph, existence or nonexistence of a feasible set
determines whether this minimum is finite or infinite.

\begin{definition}
The strong product \(G\boxtimes H\) has vertex set \(V(G)\times V(H)\).
Distinct vertices \((g,h)\) and \((g',h')\) are adjacent exactly when
each coordinate is either equal or adjacent in its factor:
\[
 (g=g'\text{ or }gg'\in E(G))
 \quad\text{and}\quad
 (h=h'\text{ or }hh'\in E(H)).
\]
\end{definition}
The distinctness condition excludes loops. Equivalently, a product
edge moves in the first coordinate, in the second, or in both
simultaneously.

\begin{lemma}\label{lem:product-distance}
For connected graphs \(G,H\),
\[
 \dist_{G\boxtimes H}((g,h),(g',h'))
   =\max\{\dist_G(g,g'),\dist_H(h,h')\}.
\]
\end{lemma}
\begin{proof}
Project a product walk of length \(\ell\) to either factor and remove
steps during which that coordinate stays fixed. The resulting factor
walk has length at most \(\ell\), proving the lower bound.

Conversely, take shortest factor paths of lengths \(r,s\).
Move along both at each step until the shorter path ends, then keep
its endpoint fixed while completing the other path. Every step is
an edge of the strong product, and the resulting walk has length
\(\max(r,s)\).
\end{proof}

Let \(V=\F_2^6\), with coordinatewise addition. For \(x\in V\), its
Hamming weight \(\wt(x)\) is the number of nonzero coordinates, and
\[
 \delta(x,y)=\wt(x+y)
\]
is the Hamming distance. Simultaneous translation and permutation of
coordinates preserve \(\delta\). We write \(0\) for the zero word,
\(\mathbf1=(1,1,1,1,1,1)\), and \(\bar a=a+\mathbf1\).
The graph \(Q_6\) joins exactly the pairs at Hamming distance one.

\begin{lemma}\label{lem:cube-distance}
The graph \(Q_6\) is connected and
\(\dist_{Q_6}(x,y)=\delta(x,y)\).
\end{lemma}
\begin{proof}
Each edge flips one coordinate, so any walk from \(x\) to \(y\) must
flip each differing coordinate at least once. Its length is at least
\(\delta(x,y)\). Flipping precisely those coordinates gives a path
of that length.
\end{proof}

\section{The obstruction in the single cube}\label{sec:factor}
The following proof applies to arbitrary subsets. No linearity,
affine-subspace structure, prescribed cardinality or parity restriction
is imposed.

\begin{lemma}\label{lem:plotkin}
If distinct words of \(D\subseteq V\) are at Hamming distance at least
four, then \(|D|\le4\). If \(|D|=3\) or \(4\), every distance between
distinct members of \(D\) is exactly four.
\end{lemma}
\begin{proof}
For \(m\) binary words \(x_1,\ldots,x_m\), let \(t_j\) be the number
whose \(j\)-th coordinate is one. Counting pairs that differ in
each coordinate gives
\begin{equation}\label{eq:distance-sum}
 \sum_{1\le i<k\le m}\delta(x_i,x_k)
    =\sum_{j=1}^6 t_j(m-t_j).
\end{equation}
If five words have pairwise distance at least four, the left side is
at least \(4\binom52=40\). On the right each contribution is at most
\(2\cdot3=6\), so the sum is at most 36, a contradiction. Every larger
set contains five words.

For three words the upper bound from~\eqref{eq:distance-sum} is
\(6\cdot2=12\), exactly the lower bound \(4\binom32\).
For four words the bounds are both
\(6\cdot4=24=4\binom42\). Equality of the sum and the pairwise lower
bounds forces every pair distance to be four.
\end{proof}

\begin{lemma}\label{lem:zero-pairs}
Translate a three- or four-word distance-four packing so that it
contains \(0\). Every other word has weight four, and its two zero
coordinates form a pair. These zero pairs are mutually disjoint.
For four centers, the three zero pairs partition the six coordinates.
\end{lemma}
\begin{proof}
The first assertion follows from Lemma~\ref{lem:plotkin}. If \(b,c\)
have weight four and zero sets \(Z_b,Z_c\) of size two, then
\[
 \delta(b,c)=|Z_b\mathbin{\triangle} Z_c|
           =4-2|Z_b\cap Z_c|.
\]
This equals four precisely when the zero sets are disjoint. Three
disjoint two-element sets exhaust the six coordinates.
\end{proof}

\begin{lemma}\label{lem:ball-intersections}
A radius-two ball in \(Q_6\) has 22 vertices. The radius-two balls
about two distance-four centers intersect in six vertices. For three
pairwise distance-four centers their three balls have a unique common
vertex. For four such centers the fourfold intersection is empty.
\end{lemma}
\begin{proof}
The first count is \(1+6+\binom62=22\). For two centers translate one
to \(0\) and call the other \(b\), with \(\wt(b)=4\). If \(x\) is in
both balls, then
\[
 4=\delta(0,b)\le\delta(0,x)+\delta(x,b)\le4.
\]
Equality forces both distances to be two and \(x\) to be supported
on two of the four nonzero coordinates of \(b\). Conversely every
such \(x\) is in the intersection. There are \(\binom42=6\).

For three centers, translate them to \(0,b,c\). Their two disjoint
zero pairs leave two coordinates \(R\) outside \(Z_b\cup Z_c\).
A word in all three balls must be a weight-two word supported on
the intersection of the supports of \(b,c\), which is precisely \(R\).
Thus the word with support \(R\) is the unique common vertex.

For four centers \(0,b,c,e\), the third zero pair is \(Z_e=R\).
The word supported on \(R\) differs from \(e\) in all six coordinates,
so it is not in the fourth ball. The same argument applies after
choosing any of the centers as the translated zero; every triple
therefore has exactly one common vertex.
\end{proof}

\begin{proposition}\label{prop:factor}
Every Hamming-distance-four packing \(D\subseteq V\) has a vertex
at distance greater than two from every member of \(D\).
Consequently \(\gamma_2^3(Q_6)=\infty\).
\end{proposition}
\begin{proof}
Lemma~\ref{lem:plotkin} gives at most four members. With zero, one
or two centers, the union of their radius-two balls has size at most
44, less than 64.

For three centers Lemma~\ref{lem:ball-intersections} and
inclusion--exclusion give
\[
 \left|\bigcup_{a\in D}B_2^{Q_6}(a)\right|
       =3\cdot22-3\cdot6+1=49.
\]
For four centers the corresponding count is
\[
 \left|\bigcup_{a\in D}B_2^{Q_6}(a)\right|
       =4\cdot22-6\cdot6+4\cdot1=56.
\]
Neither union covers \(V\). Lemma~\ref{lem:cube-distance} identifies
these Hamming balls with graph balls and proves the obstruction.
\end{proof}

\begin{remark}\label{rem:holes}
There is also a direct witness for three or four centers. Complete
the zero pairs of Lemma~\ref{lem:zero-pairs} to a partition into
three pairs. A word choosing exactly one coordinate from each pair
is at distance three from \(0\) and from every weight-four center
in the packing. The eight such words are all omitted.
This explains the obstruction without the union counts, though
the counts specify the total covered size.
\end{remark}

\section{A graph built from compatible supports}\label{sec:construction}
Define the set of unordered distance-four label pairs by
\[
 \pairs=\{\{a,b\}\subseteq V:\delta(a,b)=4\}.
\]
Every pair has two distinct members. Each label has \(\binom64=15\)
labels at distance four, so
\begin{equation}\label{eq:pair-count}
 |\pairs|=\frac{64\cdot15}{2}=480.
\end{equation}

\begin{definition}\label{def:H}
The vertices of \(H\) form the disjoint tagged union
\[
 \{A_a:a\in V\}
 \ \sqcup\ \{S_a:a\in V\}
 \ \sqcup\ \{E_P:P\in\pairs\}.
\]
We call \(A_a\) a code vertex. The other vertices are called Steiner
vertices and have supports
\[
 \sigma(S_a)=\{a\},\qquad \sigma(E_P)=P.
\]
The edges of \(H\) are precisely these unordered pairs:
\begin{enumerate}
\item \(A_aT\), where \(T\) is a Steiner vertex and \(a\in\sigma(T)\);
\item \(TU\), where \(T,U\) are distinct Steiner vertices and every
      pair of distinct words in \(\sigma(T)\cup\sigma(U)\)
      has Hamming distance at least four.
\end{enumerate}
There are no code--code edges, no loops and no other edges.
\end{definition}

The support union in the second clause is a set. A shared word is
counted once and is never compared with itself. In particular \(S_a\)
may be adjacent to \(E_{\{a,b\}}\). Tags distinguish code vertices
from singleton Steiner vertices despite their shared label.
Adjacency is symmetric and excludes equality, defining an ordinary
simple undirected graph.

\begin{proposition}\label{prop:H-connected}
The graph \(H\) has 608 vertices and is connected.
\end{proposition}
\begin{proof}
Equation~\eqref{eq:pair-count} gives \(64+64+480=608\) vertices.
For a coordinate word \(e_i\), there is a two-step path
\[
 S_a-S_{\bar a}-S_{a+e_i},
\]
because the consecutive label distances are six and five.
Thus every single-coordinate flip is realizable between singleton
vertices. Repeating such paths connects all singleton vertices.
Each \(A_a\) is adjacent to \(S_a\), and each \(E_{\{a,b\}}\) is
adjacent to \(A_a,A_b\). Every vertex belongs to the same component.
\end{proof}

\begin{lemma}\label{lem:short-path}
If \(a\ne b\) and \(\dist_H(A_a,A_b)\le3\), then \(\delta(a,b)\ge4\).
\end{lemma}
\begin{proof}
A path of length one is impossible because no code--code edge exists.
A path of length two has the form \(A_a-T-A_b\), with both labels
in \(\sigma(T)\). A singleton cannot support two distinct labels,
so \(T=E_{\{a,b\}}\) and \(\delta(a,b)=4\).

A path of length three has the form \(A_a-T-U-A_b\) with both
internal vertices Steiner: an internal code vertex would create
a code--code edge at one of the ends. The edge \(TU\) makes distinct
words in their support union separated by at least four. The union
contains \(a,b\), giving the conclusion. These cases include
shared-endpoint supports and exhaust all paths of positive length
at most three.
\end{proof}

Define the candidate centers by
\begin{equation}\label{eq:centers}
 \code=\{(a,A_a):a\in V\}.
\end{equation}
There are exactly 64 distinct product vertices in this set.

\begin{proposition}\label{prop:packing}
The set \(\code\) is a \(3\)-packing in \(Q_6\boxtimes H\).
\end{proposition}
\begin{proof}
Consider distinct labels \(a,b\). If \(\delta(a,b)\ge4\), the first
coordinate separates their centers by at least four. Otherwise
Lemma~\ref{lem:short-path} and connectedness give
\(\dist_H(A_a,A_b)\ge4\). Lemma~\ref{lem:product-distance}
proves separation at least four in either case.
\end{proof}

\section{Explicit covering sets of labels}\label{sec:labels}
For \(h\in V(H)\), put
\[
 L(h)=\{a\in V:\dist_H(h,A_a)\le2\}.
\]
Product coverage follows if each \(L(h)\) radius-two covers the cube.
We first establish the two label facts needed for this claim.

\begin{lemma}\label{lem:single-cover}
For every \(a\in V\), the set
\[
 L_a=\{a\}\cup\{b\in V:\delta(a,b)=4\}
\]
radius-two covers \(V\).
\end{lemma}
\begin{proof}
Translate \(a\) to zero. A word of weight zero, one or two is
covered by zero. A weight-three word extends to a weight-four
word at distance one. A weight-four word is itself a center.
A weight-five word contains a weight-four word at distance one,
and a weight-six word is at distance two from any weight-four
word. These cases account for every word in \(V\).
\end{proof}

\begin{lemma}\label{lem:pair-cover}
For \(a,b\in V\) with \(\delta(a,b)=4\), define
\[
 K(a,b)=\{a,b\}\cup
   \{c\in V:\delta(c,a)\ge4,\ \delta(c,b)\ge4\}.
\]
This is an eight-element set, and it radius-two covers \(V\).
\end{lemma}
\begin{proof}
Translate \(a\) to zero and permute coordinates so \(b=(1111,00)\).
Write \(c=(r,s')\), with four-bit first block \(r\) and two-bit
second block \(s'\); put \(t=\wt(r)\), \(s=\wt(s')\).
The two common-far-label inequalities become
\[
 t+s\ge4,\qquad 4-t+s\ge4.
\]
They imply \(s\ge t\) and \(s\ge4-t\).
Since \(s\le2\), necessarily \(t=s=2\).
Conversely these block weights satisfy both inequalities. Hence
\begin{equation}\label{eq:K-normalized}
 K(0,b)=\{(0000,00),(1111,00)\}
       \cup\{(u,11):u\in\F_2^4,\ \wt(u)=2\}.
\end{equation}
The six latter words are distinct from the endpoints.

For an arbitrary target \(x=(r,s')\) with block weights \(t,s\),
the endpoint distances are \(t+s\) and \(4-t+s\).
Among the six other labels the minimum distance is
\begin{equation}\label{eq:pair-minimum}
 \min_{\wt(u)=2}\delta(x,(u,11))=|t-2|+2-s.
\end{equation}
The overlap of \(\supp(u)\) with \(\supp(r)\) can be \(\min(t,2)\),
giving first-block distance
\(t+2-2\min(t,2)=|t-2|\). Such a two-element support exists for
every \(0\le t\le4\), and the second-block distance to \(11\) is \(2-s\).

It remains to check
\begin{equation}\label{eq:three-minimum}
 \min\{t+s,\ 4-t+s,\ |t-2|+2-s\}\le2.
\end{equation}
For \(t=0,4\), use the appropriate endpoint.
For \(t=1,3\), use an endpoint when \(s\le1\), and the third
expression when \(s=2\). For \(t=2\), the third expression is
\(2-s\le2\). Thus every target is covered.
\end{proof}

The minimum in~\eqref{eq:three-minimum} is displayed below.
The proof realizes each value, so this table is exact.
\[
\begin{array}{c|ccc}
 t\backslash s&0&1&2\\ \hline
 0&0&1&2\\
 1&1&2&1\\
 2&2&1&0\\
 3&1&2&1\\
 4&0&1&2
\end{array}
\]

\begin{proposition}\label{prop:label-cover}
For every auxiliary vertex \(h\), the set \(L(h)\) radius-two covers \(V\).
\end{proposition}
\begin{proof}
We treat all three types in Definition~\ref{def:H}.
For \(h=A_a\), the label \(a\) belongs to \(L(h)\).
Each \(b\) at distance four from \(a\) also belongs, via
\[
 A_a-E_{\{a,b\}}-A_b.
\]
Thus \(L_a\subseteq L(A_a)\), and
Lemma~\ref{lem:single-cover} proves the claim.

For \(h=S_a\), the edge \(S_aA_a\) reaches label \(a\).
For every \(b\) at distance four from \(a\), the union
\(\{a,b\}\) is compatible, giving
\[
 S_a-S_b-A_b.
\]
Consequently \(L_a\subseteq L(S_a)\), and the same lemma applies.

For \(h=E_{\{a,b\}}\), both endpoint code vertices are adjacent to \(h\).
A common far label \(c\) is distinct from \(a,b\). The support
union \(\{a,b,c\}\) is separated by at least four, giving
\[
 E_{\{a,b\}}-S_c-A_c.
\]
Therefore \(K(a,b)\subseteq L(E_{\{a,b\}})\).
Lemma~\ref{lem:pair-cover} completes the final case.
\end{proof}

\section{The product witness and the main theorem}\label{sec:main-proof}
\begin{proposition}\label{prop:domination}
The set \(\code\) of~\eqref{eq:centers} is a \(2\)-distance
dominating set in \(Q_6\boxtimes H\).
\end{proposition}
\begin{proof}
Fix an arbitrary product vertex \((x,h)\).
Proposition~\ref{prop:label-cover} supplies one label \(a\in L(h)\)
with \(\delta(x,a)\le2\). For that same label,
\(\dist_H(h,A_a)\le2\). Hence
\[
 \dist_{Q_6\boxtimes H}((x,h),(a,A_a))
    =\max\{\delta(x,a),\dist_H(h,A_a)\}\le2.
\]
The center \((a,A_a)\) belongs to \(\code\).
The argument covers every \(x\) and every type of \(h\),
thus all \(64\cdot608=38,912\) product vertices.
\end{proof}

\begin{proof}[Proof of Theorem~\ref{thm:main}]
Proposition~\ref{prop:factor} proves \(\gamma_2^3(Q_6)=\infty\).
Proposition~\ref{prop:H-connected} proves that \(H\) is connected
and has 608 vertices. Propositions~\ref{prop:packing} and
\ref{prop:domination} give the feasible set \(\code\) of size 64.
Thus the product packing-domination number is finite and at most 64.
The edge count is proved in Proposition~\ref{prop:edge-count} below.
\end{proof}

\begin{remark}\label{rem:projection}
The first-coordinate projection of \(\code\) is the whole cube,
which is not a distance-four packing. This is compatible with
product packing because the second coordinate separates every
incompatible pair of labels. Its radius-two neighborhoods still
contain enough labels to cover the cube. Projecting a product
witness to one factor therefore does not preserve packing.
\end{remark}

\section{Edge count and the close-vertex condition}\label{sec:additional}
\begin{proposition}\label{prop:edge-count}
The graph \(H\) has exactly 9,168 edges.
\end{proposition}
\begin{proof}
Partition the edges by vertex type. There are 64 code--singleton
incidences and \(2\cdot480\) code--pair incidences, giving 1,024.
For singleton--singleton edges a fixed label has
\(\binom64+\binom65+\binom66=22\) compatible distinct labels.
Division by two gives \(64\cdot22/2=704\).

A fixed pair \(\{a,b\}\) is compatible with singleton supports
at \(a,b\) and at the six common far labels of
Lemma~\ref{lem:pair-cover}. These are exactly the compatible
singleton supports, giving \(480\cdot8=3,840\) such edges.

For pair--pair edges sharing an endpoint, translate it to zero.
The other two labels have weight four, and compatibility means
their zero pairs are disjoint. There are 15 choices of the first
zero pair and \(\binom42=6\) disjoint choices of the second.
Division by two gives 45 edges at each shared endpoint, hence
\(64\cdot45=2,880\). The shared endpoint is unique.

For disjoint pair supports, their union is a four-word
distance-four packing. By Lemma~\ref{lem:zero-pairs}, a packing
containing a specified label is determined by a partition into
three unordered zero pairs. There are
\[
 \frac{6!}{(2!)^3\,3!}=15
\]
such partitions. Conversely, each partition gives three
weight-four words whose zero pairs are its blocks; together
with zero they form a valid four-word distance-four packing.
Choosing a specified label in 64 ways counts
every packing four times, giving \(64\cdot15/4=240\) packings.
Each has three partitions into two unordered pair supports,
giving \(240\cdot3=720\) edges. The total is
\[
 1,024+704+3,840+2,880+720=9,168.
\]
\end{proof}

Theorem~3.3 of~\cite{bickz-strong} proves the conjectured
implication when the \emph{other} factor \(H\) has a
\((d,p)\)-close vertex: a vertex \(h\) such that all distances
between vertices in \(B_d^H(h)\) are at most \(p\).
The next observation checks consistency with that result.

\begin{proposition}\label{prop:no-close}
Every radius-two ball of \(H\) has diameter at least four.
In particular \(H\) has no \((2,3)\)-close vertex.
\end{proposition}
\begin{proof}
For \(h=A_a\) or \(S_a\), choose two labels in \(L_a\)
whose translated weight-four supports share three coordinates,
for example \(\{1,2,3,4\}\) and \(\{1,2,3,5\}\).
The labels have Hamming distance two and both code vertices
lie in \(B_2^H(h)\).

For \(h=E_{\{a,b\}}\), normalize as in
Lemma~\ref{lem:pair-cover}. Choose the labels
\((u,11),(v,11)\) with first-block supports
\(\{1,2\},\{1,3\}\). Again their Hamming distance is two
and both code vertices belong to \(B_2^H(h)\).

Lemma~\ref{lem:short-path} puts their \(H\)-distance at least
four, while their paths through \(h\) put it at most four.
Thus each ball contains a pair at distance exactly four.
\end{proof}

\section{Formal verification}\label{sec:formal}
The complete existence assertion has been checked in
Lean~\cite{lean4} using Mathlib~\cite{mathlib}.
The public sources and source-specific receipts accompany
this paper~\cite{repository}. The environment is pinned to
Lean 4.34.1 and Mathlib commit
\[
 \texttt{d13f23b723b8a846827a245b89c10fc7d3f11612}.
\]
The dependency manifest is supplied in the isolated package
used by the replay script.

\subsection{Objects and the path bridge}
Cube vertices have type \texttt{BitVec 6}. Hamming distance
counts all six differing bits and returns a natural number;
there is no bit-vector overflow in the distance result.
The actual \texttt{SimpleGraph} \(G\) has adjacency
\(\delta(a,b)=1\). Auxiliary vertices form a tagged sum
of 64 code labels, 64 singleton labels and canonical
unordered distance-four pairs. Strict numerical ordering
selects one representative of every unordered pair.
The source defines the symmetric loopless adjacency of \(H\)
using support membership and compatibility. It proves its
product adjacency equivalent to the three ordinary
strong-product edge cases. No abstract distance table
is assumed.

For convenient bounded lengths the formalization uses padded
reachability. Write \(\operatorname{Step}_F(u,v)\) for equality
or adjacency in \(F\), and define
\begin{align*}
 \Reach_F(0,u,v)&\ \Longleftrightarrow\ u=v,\\
 \Reach_F(n+1,u,v)&\ \Longleftrightarrow\
   \exists w,\ \operatorname{Step}_F(u,w)\land\Reach_F(n,w,v).
\end{align*}
The theorem \texttt{reach\_iff\_walk} proves equivalence with
existence of a standard Mathlib \(F\)-walk of length at most \(n\).
Equality steps pad a shorter walk; they do not add graph loops.
The proof converts padded reachability to a walk and pads a
shorter walk in the converse direction.

The theorem \texttt{G\_reach\_iff\_hamming} proves
\[
 \Reach_G(n,a,b)\ \Longleftrightarrow\ \delta(a,b)\le n
\]
for every \(n\). The product bridge \texttt{strong\_reach} proves
\[
 \Reach_{G\boxtimes H}(n,(x,h),(a,k))
 \ \Longleftrightarrow\
 \Reach_G(n,x,a)\land\Reach_H(n,h,k).
\]
The same endpoint \((a,k)\) appears throughout, matching the
usual strong-product distance formula.

\subsection{The formal endpoint and its interpretation}
The definition \texttt{WalkPacking F S} states that no standard
walk of length at most three joins distinct members of \(S\).
The definition \texttt{WalkCovers2 F S} states that every vertex
has a standard walk of length at most two to some member of \(S\).
On connected graphs these are exactly our packing and domination
properties.

The declaration \texttt{original\_existence\_counterexample}
has the mathematical content
\begin{align}
 &\neg\exists S\subseteq V(G),\
        \operatorname{WalkPacking}(G,S)
        \land\operatorname{WalkCovers2}(G,S),
       \label{eq:formal-factor}\\
 &\exists S\subseteq V(G)\times V(H),\
        \operatorname{WalkPacking}(G\boxtimes H,S)
        \land\operatorname{WalkCovers2}(G\boxtimes H,S).
       \label{eq:formal-product}
\end{align}
The second witness is the defined set \(\code\).
The first quantifies over every cube-vertex set, not just linear
codes or finite samples. The endpoint has no unproved geometric
premise, external-certificate hypothesis or extra center assumption.

On finite graphs, these statements are the existence content
of infinite and finite packing-domination numbers. A separate
numerical function \(\gamma_d^p\) is not defined in the Lean
sources. The written proof gives the explicit cardinality
\(|\code|=64\), hence the bound in Theorem~\ref{thm:main}.
No numerical minimization or optimality claim is imported
into the formal endpoint.

\subsection{Proof structure, audits and remaining boundaries}
There are three project modules. The existing factor module
proves a missing-vertex theorem for arbitrary subsets of
six-bit Hamming space. It normalizes a nonempty set by XOR
translation, lists the 22 possible nonzero labels at distance
at least four from zero, and constructs exact missing-vertex
witnesses. Ordinary \texttt{decide} proves finite certificate
lemmas, including a \(22^3\) compatibility certificate.
No power set of 64 vertices is enumerated. This is an
alternative formal proof of Proposition~\ref{prop:factor},
rather than a formalization of its inclusion--exclusion text.

The new geometry module proves Hamming metric facts, a one-bit
distance-reducing step and both label-covering facts. It uses
XOR normalization and ordinary kernel-checked \texttt{decide}
on small closed finite statements. The graph module supplies
actual graphs, bounded-walk bridges, short-center exclusion,
same-center coverage and the final conjunction.

Primary records contain 12 printed axiom audits for the geometry
module and 26 for the graph module. A separately run replay
builds the factor, geometry and graph modules from their actual
sources in a fresh project object directory, in dependency order.
It records 12, 12 and 26 audits respectively. All recorded runs
have exit code zero, no warnings or errors and no nonstandard
axioms. Audit records inspect declarations; they are not a
count of newly discovered mathematical theorems.

The final endpoint uses only the standard Lean axioms
\texttt{propext}, \texttt{Classical.choice} and \texttt{Quot.sound}.
The sources use no \texttt{sorry}, \texttt{native\_decide}, added
axiom, solver oracle or Python certificate as a proof.
Public records retain source hashes, UTC intervals, logs and
every printed dependency list.

The cardinality 608 of the auxiliary vertex type is a Lean
theorem. Connectedness of \(H\), its 9,168-edge count and its
exact diameter are not separate Lean endpoints in this version.
Connectedness and the edge count have complete written proofs
above and exact graph-program checks. A numerical
\texttt{SimpleGraph.dist} definition is not used; ordinary
shortest paths are represented by standard walks and lengths.
These boundaries do not weaken the formal existence assertion,
but they must not be hidden by claiming that every auxiliary
statistic has been formalized.

\section{Exact computational verification}\label{sec:computational}
Two separately implemented graph checkers use Python's standard
library and exact integer and set operations. Neither uses a
floating-point optimization or solver. The mathematical proof
of Theorem~\ref{thm:main} does not depend on their output.

\subsection{First construction and witness certificate}
The first checker represents a cube word as an integer from
0 to 63 and computes Hamming distance by XOR and population
count. It constructs cube edges and verifies that breadth-first
search distances equal Hamming distances for all cube vertices.
It constructs all 608 tagged auxiliary vertices and every edge
of Definition~\ref{def:H}.

Breadth-first searches from all 64 code vertices give distances
to every auxiliary vertex. All \(\binom{64}{2}=2,016\) center
pairs are tested using the maximum factor distance.
For each of the 38,912 product targets the program saves one
label satisfying both distance bounds and rechecks both bounds
for every saved same-center witness. The public graph and
certificate files contain the tagged adjacency data and the
distance/witness arrays.

\subsection{Separate reconstruction and shortest-distance checks}
The second checker initially reads no graph or certificate from
the first implementation. It reconstructs the vertices and
support-union adjacency directly from the definitions, then
compares adjacency with a separately expressed predicate on
all \(\binom{608}{2}=184,528\) unordered distinct vertex pairs.
It checks for missing and spurious edges.

It runs breadth-first search from all 608 vertices, obtaining
369,664 distances. These also undergo a shortest-distance
certificate check: for every source \(v\), the recorded function
\(f_v\) is zero at \(v\), positive elsewhere, differs by at most
one across every edge, and each non-source vertex \(w\) has a
neighbor of value \(f_v(w)-1\). The edge bound makes every
\(v\)-to-\(w\) walk at least \(f_v(w)\) steps long; successive
descending neighbors construct a walk of exactly that length.
These conditions certify shortest distances, rather than merely
the shape of a BFS array. Symmetry is checked too.

The checker independently tests all 2,016 center pairs and
38,912 same-center targets. It verifies that the sufficient
label sets of Section~\ref{sec:labels} lie in the actual
radius-two neighborhoods and cover the cube. Only after passing
this phase does it read the first checker's artifacts.
The comparison matches all 608 canonical vertices, 9,168 edges,
38,912 code-distance entries and all covering witnesses.

The recorded independent run used Python 3.12.14 on 64-bit
Windows. The scripts require Python 3.10 or later.
There are no recorded uncovered targets or packing violations,
and the minimum distance between the specified product centers
is four. The complete BFS also records diameter four.
That statistic is supplementary: the proof needs only
connectedness, short-path exclusion and radius-two coverage.

\subsection{Meaning of separate checks}
Separate implementations, reconstruction before reading the
first artifacts and a fresh Lean object directory reduce
some common checking errors. They do not establish independent
human peer review, originality or journal acceptance.
AI agents participated in implementation and review.
Exact-program outputs are finite evidence, the written
argument is a complete mathematical proof, and the Lean
endpoint is a kernel-checked existence assertion.
These are related but distinct forms of evidence.

\section{Discussion and limits}\label{sec:discussion}
The obstruction in \(Q_6\) is not inherited by every product
containing it. The auxiliary graph supplies a separation
coordinate while retaining the requirement that a single
center dominate both coordinates. At a pair Steiner vertex,
the endpoints alone do not cover the cube. The six additional
labels in~\eqref{eq:K-normalized} fill the missing coverage while
remaining compatible with the pair support.
This eight-label calculation enables the global construction.

The theorem refutes the exact universal implication of
Conjecture~3.1 in~\cite{bickz-strong}. It does not contradict
perfect-code inheritance at \(p=2d\)~\cite{perfect}: there the
balls are disjoint and each target has a unique dominating
center. Our balls may overlap. Nor does it contradict the
close-vertex theorem; Proposition~\ref{prop:no-close} shows
that its \(H\)-side hypothesis is absent.
The complexity, cycle and tree results
of~\cite{bickz-complexity} concern other assertions.

We do not determine \(\gamma_2^3(Q_6\boxtimes H)\) exactly,
minimize the number of auxiliary vertices or classify all
factors repairing the cube's feasibility obstruction.
The 64 centers are a witness, not an asserted optimum.
Replacing \(H\) with another copy of \(Q_6\) is a different
problem: this paper does not decide whether
\(\gamma_2^3(Q_6\boxtimes Q_6)\) is finite.

The primary preprint and published article were checked on
9 and 10 October 2026; their conjecture statements agree in the
relevant scope. Targeted identifier, author, older
\((p,d)\)-domination and support-compatibility searches did
not locate a verified earlier complete result subsuming this
construction. Those searches are not exhaustive and do not
establish firstness. Earlier product bounds such
as~\cite{fisher} also involve packing and domination.
A full historical audit remains necessary before asserting
priority. Correctness of the explicit counterexample and
historical priority are separate questions.
The accompanying publication audit records the checked sources,
the limits of the citation search and the precise formalization
boundary. It is not a bibliometric similarity certificate or an
external referee report. Before submission, the eventual human
authors must examine the proof, the cited literature and the
accountability requirements of the selected venue.

\section{AI involvement and academic status}\label{sec:disclosure}
Substantial generative-AI assistance was used in this research.
The work used OpenAI Codex with a GPT-6-based agent system;
no more specific model version is asserted here.
AI systems participated in proposing the support-based graph,
developing and checking derivations, writing Lean sources,
implementing exact graph checkers, inspecting statements and
axiom lists, conducting targeted literature searches, and
drafting and revising this manuscript. Separate AI-agent
mathematical reviews and program reconstructions were performed;
they are not described as independent human reviews.
The Lean kernel checks and actual program executions are
recorded separately from AI-generated reasoning.

As of 10 October 2026, the project owner has supplied the draft human
author name River Zhang, ORCID identifier 0009-0004-2437-8566,
and affiliation Chengdu Neusoft University. Accountable human examination,
license selection and external publication approval remain pending.
No AI system is listed as an author. The human author must examine and take
responsibility for all content, including proofs, references,
verification claims and disclosures.
This public research draft does not claim completed human
scholarly supervision or peer review. It has not been submitted
to arXiv or a journal, has no arXiv identifier of its own, and
does not claim acceptance or compliance with a particular
venue's eligibility requirements.
No Zenodo record or DOI has been created for this manuscript by
this preparation work. The deposit license has not been selected.
The author information is
recorded from the project owner's instruction, not invented by an AI
system. The substantial mathematical and writing assistance described
above must not be represented as language polishing alone.
No correspondence address, department, funding declaration or
statement of the human author's individual research contributions
has been supplied. These items must be confirmed by the author before
submission; no absence of funding or conflicts of interest is inferred.

\appendix
\section{Source correspondence and reproducibility}\label{app:repro}
The repository~\cite{repository} supplies the following material.
The factor path is relative to the repository root. Other paths in
the table are relative to
\path{006-strong-product-packing-counterexample/}.
They identify public source files rather than a private machine environment.

\begin{longtable}{@{}>{\raggedright\arraybackslash}p{0.29\textwidth}>{\raggedright\arraybackslash}p{0.65\textwidth}@{}}
\toprule
Mathematical content & Public source and verification boundary\\
\midrule
\endhead
Arbitrary-set cube obstruction &
\path{notes/proof/Q6PackingDomination.lean}.
The missing-vertex theorem covers every subset of
\texttt{BitVec 6}.\\[5pt]
Hamming geometry and label coverage &
\path{proof/Q6PackingComplexGeometry.lean}.
The declarations \texttt{single\_cover} and \texttt{pair\_cover}
give actual translated covering witnesses.\\[5pt]
Graphs, walks and complete endpoint &
\path{proof/Q6PackingComplexCounterexample.lean}.
Its \texttt{original\_existence\_counterexample} is
the conjunction~\eqref{eq:formal-factor}--\eqref{eq:formal-product}.\\[5pt]
Compilation and axiom evidence &
\path{results/lean-verification.json}
and source-specific logs. The fresh rebuild compiles all three
project sources.\\[5pt]
Explicit graph and witness &
\path{results/graph.json}
and \path{results/certificate.json}.\\[5pt]
Exact reconstruction and comparison &
The scripts \texttt{verify\_graph.py} and
\texttt{verify\_independent.py} in the topic directory,
with graph and comparison receipts under \texttt{results/}.\\[5pt]
Connectedness and 9,168 edges &
Propositions~\ref{prop:H-connected} and~\ref{prop:edge-count},
plus exact graph replay; not separate Lean theorems here.\\[5pt]
Numerical bound and limitations &
The injective label map gives \(|\code|=64\) in the written proof.
The Lean endpoint defines no numerical \(\gamma_d^p\) function
and supplies no optimality theorem.\\
\bottomrule
\end{longtable}

For a fresh replay, install the pinned Lean toolchain and
prepare the dependency package
\path{002-weighted-rectangular-pruning/proof/mathlib}.
Its manifest fixes Mathlib and the transitive dependencies.
Obtain the required dependency artifacts using Lake's cache
command if they are absent. In a fresh checkout, first run
\texttt{lake update} and \texttt{lake exe cache get} from that
dependency package; its Git requirement fixes the Mathlib revision.
The versioned supplementary archive includes the toolchain,
Git dependency configuration, transitive manifest and the external
factor module, rather than requiring a private local dependency path.
The replay first queries the actual Lean version and Mathlib Git
revision and checks them against the pins. An environment-only
check does not count as a proof compilation. The replay runs all three
project sources serially and creates new object files;
it does not accept old project objects as a substitute
for source compilation.

Run the following commands from the repository root, with
new output directories. Relative paths below are examples
usable on Windows or Unix-like systems.
\begin{verbatim}
python 006-strong-product-packing-counterexample/verify_graph.py \
  q6-primary-new
python 006-strong-product-packing-counterexample/verify_independent.py \
  own --out q6-independent-new
python 006-strong-product-packing-counterexample/verify_independent.py \
  compare --out q6-independent-new --root q6-primary-new
python 006-strong-product-packing-counterexample/verify_lean.py \
  --out q6-lean-new
\end{verbatim}
The displayed backslash continuations are Unix-like shell syntax.
In PowerShell put each command on one line or use that shell's
continuation syntax. The option \texttt{--lake} accepts an
explicit Lake executable path.

The receipts retain source and artifact hashes and distinguish
reconstruction without the first checker's data from subsequent
comparison. A regenerated result should be assessed against
the actual source revision and precise checked assertions.
An old result file or matching timestamp is not a fresh execution.

\begin{thebibliography}{10}
\raggedright
\bibitem{beineke-henning}
L.~W.~Beineke and M.~A.~Henning,
\emph{Some extremal results on independent distance domination in graphs},
Ars Combinatoria \textbf{37} (1994), 223--233.
\url{https://combinatorialpress.com/article/ars/Volume%20037/volume-37-paper-20.pdf}.

\bibitem{bickz-strong}
C.~Bujt\'as, V.~Ir\v si\v c Chenoweth, S.~Klav\v zar and G.~Zhang,
\emph{On \(d\)-distance \(p\)-packing domination number in strong products},
Annals of Combinatorics (2026), published 26 March 2026.
\href{https://doi.org/10.1007/s00026-026-00814-0}
{doi:10.1007/s00026-026-00814-0}.
Original preprint:
\href{https://arxiv.org/abs/2510.02749v1}{arXiv:2510.02749v1},
3 October 2025.

\bibitem{bickz-complexity}
C.~Bujt\'as, V.~Ir\v si\v c Chenoweth, S.~Klav\v zar and G.~Zhang,
\emph{The \(d\)-distance \(p\)-packing domination number:
complexity, cycles, and trees},
Aequationes Mathematicae \textbf{100} (2026), article 25.
\href{https://doi.org/10.1007/s00010-026-01266-w}
{doi:10.1007/s00010-026-01266-w}.

\bibitem{perfect}
G.~Abay-Asmerom, R.~H.~Hammack and D.~T.~Taylor,
\emph{Perfect \(r\)-codes in strong products of graphs},
Bulletin of the Institute of Combinatorics and its Applications
\textbf{55} (2009), 66--72.
\url{https://richardhammack.github.io/reprints/rperfect_codes_strong_prod.pdf}.

\bibitem{fisher}
D.~C.~Fisher,
\emph{Domination, fractional domination, 2-packing, and graph products},
SIAM Journal on Discrete Mathematics \textbf{7} (1994), no.~3, 493--498.
\href{https://doi.org/10.1137/S0895480191217806}
{doi:10.1137/S0895480191217806}.

\bibitem{handbook}
R.~Hammack, W.~Imrich and S.~Klav\v zar,
\emph{Handbook of Product Graphs}, second edition,
CRC Press, Boca Raton, 2011.

\bibitem{plotkin}
M.~Plotkin,
\emph{Binary codes with specified minimum distance},
IRE Transactions on Information Theory \textbf{6} (1960), no.~4, 445--450.
\href{https://doi.org/10.1109/TIT.1960.1057584}
{doi:10.1109/TIT.1960.1057584}.

\bibitem{lean4}
L.~de Moura and S.~Ullrich,
\emph{The Lean 4 theorem prover and programming language},
in \emph{Automated Deduction---CADE 28},
Lecture Notes in Computer Science \textbf{12699},
Springer, 2021, 625--635.
\href{https://doi.org/10.1007/978-3-030-79876-5_37}
{doi:10.1007/978-3-030-79876-5\_37}.

\bibitem{mathlib}
The mathlib Community,
\emph{The Lean mathematical library},
in \emph{Proceedings of the 9th ACM SIGPLAN International
Conference on Certified Programs and Proofs},
ACM, 2020, 367--381.
\href{https://doi.org/10.1145/3372885.3373824}
{doi:10.1145/3372885.3373824}.

\bibitem{repository}
Accompanying mathematical sources and verification records,
\emph{006: strong-product packing counterexample},
\url{https://github.com/mio-qwq/math/tree/613108f1fa107a88a1049ed5f25e6b9313d576ed/006-strong-product-packing-counterexample}.
Public source baseline for this draft:
commit \texttt{613108f} (9 October 2026).
\end{thebibliography}
\end{document}
